

\documentclass[11pt]{article}

\usepackage[margin=1in]{geometry}
\usepackage{amsmath,amssymb,amsthm,mathtools}
\usepackage{graphicx}
\graphicspath{{figs/}}
\usepackage{booktabs}
\usepackage{enumitem}
\usepackage{hyperref}
\usepackage{microtype}
\usepackage{listings}
\usepackage{xcolor}
% Listings defaults for JSON evidence artifacts
\lstset{basicstyle=\ttfamily\footnotesize,breaklines=true,columns=fullflexible}

\newtheorem{definition}{Definition}
\newtheorem{theorem}{Theorem}
\newtheorem{lemma}{Lemma}
\newtheorem{corollary}{Corollary}
\newtheorem{proposition}{Proposition}
\newtheorem{remark}{Remark}

\newcommand{\Adv}{\mathrm{Adv}}
\newcommand{\TV}{\mathrm{TV}}

\title{Closed-View Auditing for Stateful Anonymity:\\Odometers, Alarms, and Machine-Checkable CPPCs}
\author{}
\date{Unified archive note v0.08 (2026-02-28; archive v125)}


\begin{document}
\maketitle

\begin{abstract}
Deploying anonymous overlays and anonymous DHTs requires monitoring, debugging, and drift detection.
But monitoring itself creates a control-plane leakage channel: logs, meters, and alerts reveal hidden state
and can enable sequential identification.
We propose a practical deployment pattern: a \emph{measurement firewall} that keeps rich monitoring internal (closed view),
and releases only budgeted, sanitized summaries governed by a privacy accountant (odometer/filter).
We show that publicly publishing an alert bit (or trigger time/window) incurs an unavoidable privacy ``oracle tax'':
achieving low false-alarm and high power forces a minimum spend.
Finally we propose CPPCs (Control-Plane Privacy Cards): short, machine-checkable manifests describing
which witness channels are present, what accounting model is used, and what budget is allocated,
analogous in spirit to model cards/datasheets and SBOM workflows (see ABOM/OINL, Anonymity~C), but focused on anonymity control planes.\cite{series:abomoinl}
\end{abstract}

\paragraph{Compilation note.}
This paper is brief-by-design. A companion addendum collects extended multi-stream monitoring notes and CPPC schema excerpts (State~3A).\cite{series:state3addendum}

\paragraph{Scope.}
We isolate and formalize \emph{Closed-View Auditing for Stateful Anonymity:} as a reusable contract surface in the anonymity / Anonymous-DHT series.

\paragraph{Units.} Budgets are published in bits (use $\log_2$ and $2^{\epsilon}$); results stated in nats can be converted by dividing by $\ln 2$ (see Synthesis~8).\cite{series:notation}

\paragraph{Imports.}
We import the state-dependent anonymity viewpoint (amplification bounds, odometers/filters as accountants) from AnonDHT~State~1.\cite{series:state1}
Stop-time padding and schedule compilation are treated as black-box mechanisms elsewhere in the series and appear here only as examples of witness channels CPPCs must disclose.\cite{series:stoptimeaddendum,series:schedcompiler}
CPPCs are designed to align with the series contract-object and manifest stack (Synthesis~0; ABOM/OINL, Anonymity~C).\cite{series:contractobjects,series:abomoinl}
Trace-interface naming (exposure normal forms, ENF-ID) and threat-window declarations (TW-ID) are imported from Synthesis~3/4.\cite{series:traceifaces,series:threatwindows}

\paragraph{Exports.}

\paragraph{Claim fields (non-negotiable).}
A CPPC must record which transcript projection it is constraining (\texttt{exposure\_nf\_id}) and which repetition/threat window the bound is meant to cover (\texttt{tw\_id}). For stateful monitoring surfaces, CPPCs should also carry a digest-bound \texttt{state\_decl\_id} committing to the declared state surfaces and accounting scope (see State~1 and the worked example, Synthesis~17).
These ids are claim fields, not prose: they are carried by the contract object, by receipts/certificates, and by the ABOM that commits the supporting bundle.


We export (i) the \emph{measurement firewall} pattern for keeping rich monitoring closed-view while releasing only budgeted summaries, (ii) an explicit ``alert-oracle tax'' lower bound (Cor.~\ref{cor:alert-oracle-tax}) quantifying the unavoidable spend of publishing alarms, and (iii) a CPPC (Control-Plane Privacy Card) manifest schema for machine-checkable disclosure of witness channels, accountants, and allocated budgets, treated as a first-class contract object in the archive interface (Synthesis~0).\cite{series:contractobjects}.

\paragraph{Artifact policy.}
This archive is source-only; supporting artifacts (plots, scripts, datasets, and tooling) are available on request.\cite{series:artifactpolicy}

\section{Introduction}
Anonymity systems require monitoring: to detect congestion collapse, path bias, Sybil capture, or failures.
Yet monitoring is itself a control-plane channel.
Publishing fine-grained telemetry, alarms, or ``health'' beacons can enable sequential identification of hidden state.

\paragraph{Closed view.}
We propose a deployment pattern: keep rich telemetry and detectors \emph{inside} a measurement firewall, and release only budgeted, sanitized summaries.
The core idea is to treat monitoring as an interactive mechanism that must be accounted for.

\begin{figure}[t]
  \centering
  \fbox{\parbox{0.92\linewidth}{\centering\textbf{Figure omitted in source-only archive.}\\(Plot/diagram and scripts available on request.)}}
  \caption{Measurement firewall pattern. Rich telemetry stays internal; releases are budgeted and coarse.}
  \label{fig:fw}
\end{figure}

\paragraph{Non-redundancy.}
AnonDHT~State~1 develops state-dependent anonymity games and amplification bounds~\cite{series:state1}.
Here we focus on the deployment interface: how to monitor without creating a public witness stream, and how to document/attest the control-plane privacy contract via CPPC manifests.


\section{Closed-view architecture}
\label{sec:closed}

\paragraph{Composition stack.}
A useful mental model is a composition stack: raw telemetry feeds detectors; detectors feed internal dashboards; and only selected, budgeted releases cross the firewall.

\begin{figure}[t]
  \centering
  \fbox{\parbox{0.92\linewidth}{\centering\textbf{Figure omitted in source-only archive.}\\(Plot/diagram and scripts available on request.)}}
  \caption{Closed-view composition stack.}
\end{figure}

\paragraph{Why ``closed view'' is technically meaningful.}
If the public sees only a coarsened or subsampled view of internal events, then the observer channel can contract distinguishability.
In DP analyses, this corresponds to privacy amplification by subsampling/coarsening~\cite{balle2018subsampling}.
The firewall enforces this by policy: detectors may look at everything internally, but releases are constrained.

\paragraph{Accountant placement.}
An odometer/filter is placed at the firewall boundary.
It tracks realized privacy loss (odometer) or halts publication (filter) when a budget is exceeded~\cite{rogers2016odometers}.
We emphasize adaptive-parameter regimes: operators often increase monitoring or tighten thresholds in response to events.


\section{Accounting, alarms, and the alert-oracle tax}
\label{sec:alerts}

Public alerts are witnesses.
An ``alarm bit'' is a hypothesis test output; publishing the trigger time/window turns monitoring into a sequential test channel.
This issue is explicit in the DP sequential detection literature: releasing a detector statistic and its stopping time can compromise privacy, motivating DP variants that noise both statistics and thresholds~\cite{xie2025dpcusum}.

AnonDHT~State~1 develops the general ``stability $\Rightarrow$ distinguishing advantage'' bridge for state-dependent anonymity channels; here we apply the same logic to the simplest possible public witness: an alarm bit.


\subsection{A minimal tax statement}
The relationship between DP-style stability and hypothesis testing is classical; see, e.g., max-divergence/DP as a bound on likelihood ratios and testing power~\cite{dworkroth2014foundations,balle2020rdp_testing}.

\begin{corollary}[Alert-oracle tax]
\label{cor:alert-oracle-tax}
Let $B\in\{0,1\}$ be a \emph{public} alarm bit released by a detector, and let $\sigma\in\{0,1\}$ index two adjacent/control-plane states (e.g., ``good'' vs ``bad'').
If the detector has false-alarm probability $\Pr[B{=}1\mid\sigma{=}0]\le \alpha$ and miss probability $\Pr[B{=}0\mid\sigma{=}1]\le \beta$, then any $(\varepsilon,\delta)$ stability guarantee for the \emph{alarm channel} must satisfy
\[
1-\beta \;\le\; 2^{\varepsilon}\alpha+\delta.
\]
In particular, if $1-\beta>\delta$, then $\varepsilon \ge \log_2\!\big((1-\beta-\delta)/\alpha\big)$.
\end{corollary}

\begin{proof}
Apply the $(\varepsilon,\delta)$ inequality to the event $E=\{B=1\}$:
$\Pr[B{=}1\mid\sigma{=}1]\le 2^{\varepsilon}\Pr[B{=}1\mid\sigma{=}0]+\delta$.
Use $\Pr[B{=}1\mid\sigma{=}1]\ge 1-\beta$ and $\Pr[B{=}1\mid\sigma{=}0]\le \alpha$.
\end{proof}

Figure~\ref{fig:alarm-tax} provides an illustration.

\begin{figure}[t]
  \centering
  \fbox{\parbox{0.92\linewidth}{\centering\textbf{Figure omitted in source-only archive.}\\(Plot/diagram and scripts available on request.)}}
  \caption{Illustrative alarm tax curve: tighter $(\alpha,\beta)$ implies larger minimum $\varepsilon$.}
  \label{fig:alarm-tax}
\end{figure}

\paragraph{Concrete calibration (brief).}
Table~\ref{tab:alert-tax-calibration} gives a numeric sense of Corollary~\ref{cor:alert-oracle-tax}; tight public alarms imply nontrivial \(\varepsilon\) spending.

\begin{table}[h]
\centering
\begin{tabular}{@{}rrr@{}}
\toprule
Miss prob. $\beta$ & False-alarm $\alpha$ & Lower bound $\varepsilon \gtrsim \log_2\!((1-\beta)/\alpha)$ \\
\midrule
0.10 & 0.05  & 2.89 \\
0.10 & 0.01  & 4.50 \\
0.10 & 0.001 & 6.80 \\
0.05 & 0.05  & 2.94 \\
0.05 & 0.01  & 4.55 \\
0.05 & 0.001 & 6.86 \\
0.01 & 0.05  & 2.99 \\
0.01 & 0.01  & 4.60 \\
0.01 & 0.001 & 6.90 \\
\bottomrule
\end{tabular}
\caption{Calibration for Corollary~\ref{cor:alert-oracle-tax}: tight alarms cost nontrivial $\varepsilon$.
This table is illustrative and ignores small-$\delta$ and composition details.
}
\label{tab:alert-tax-calibration}
\end{table}



\subsection{Mitigation patterns for public reporting}
The tax in Corollary~\ref{cor:alert-oracle-tax} is a design constraint: if you insist on publishing \emph{high-quality} alerts, you are committing to spending a nontrivial privacy/stability budget on the alert channel.
Conversely, if you must keep the public budget small, you must deliberately degrade the alert channel (larger $\alpha$ and/or $\beta$).

A few practical patterns are simple and composable:
\begin{enumerate}[leftmargin=*]
  \item \textbf{Fixed schedule (no data-dependent timing).} Release summaries at predetermined times; do not publish ``as soon as'' a detector fires.
  \item \textbf{Coarsen and delay.} Map the trigger time/window $T$ to a coarse bucket $B=f(T)$ and release only after a fixed delay. Post-processing cannot worsen stability/privacy: if an internal report is $(\varepsilon,\delta)$-DP then so is any $f(\cdot)$ of it.
  \item \textbf{Delayed aggregates under continual observation.} When possible, replace per-event public alarms with aggregate releases (counts/rates) on a fixed cadence, using continual-observation DP mechanisms (binary-tree / hybrid counters).
  This trades a sharp stopping-time oracle for a smoothed stream whose disclosure properties are well understood~\cite{dwork2010continual,chan2010continual,jain2023priceContinual,cohen2024continual_lowerbounds}.
  \item \textbf{Output perturbation.} If an alert must be public, add symmetric randomized-response noise to the published bit; this raises false alarms by design and can enforce a target $(\alpha,\beta)$ curve under a fixed $\varepsilon$ budget~\cite{warner1965randomizedresponse}.
  \item \textbf{Commit-then-reveal (fixed time).} When operators want accountability without a real-time oracle, publish a daily/weekly commitment to the internal alert transcript (e.g., a hash or Pedersen-style commitment) on a fixed schedule, and reveal the plaintext transcript only after a policy delay.
  The commitment anchors integrity (prevents backdating) while the delayed reveal intentionally degrades the public channel by widening the effective alert window, aligning with Corollary~\ref{cor:alert-oracle-tax}~\cite{pedersen1991distributedprovers}.
\end{enumerate}

These patterns are intentionally generic: they can be layered on top of any closed-view detector, and they make the ``alert quality vs.\ public budget'' tradeoff explicit.

\begin{lemma}[Randomized-response degradation of a public alarm]
\label{lem:rr-alarm}
Let $B\in\{0,1\}$ be an internal alarm bit with false-alarm rate $\alpha_0=\Pr[B{=}1\mid\sigma{=}0]$ and miss rate $\beta_0=\Pr[B{=}0\mid\sigma{=}1]$.
Publish $\widetilde B := B\oplus N$ where $N\sim\mathrm{Bernoulli}(q)$ is independent symmetric noise.
Then the public false-alarm and miss rates are
\[
\alpha = (1-q)\alpha_0 + q(1-\alpha_0),\qquad
\beta  = (1-q)\beta_0  + q(1-\beta_0).
\]
In particular, increasing $q$ intentionally degrades alert quality (driving $(\alpha,\beta)$ toward $(1/2,1/2)$), which can be used to satisfy a fixed public budget implied by Corollary~\ref{cor:alert-oracle-tax}.
\end{lemma}

\subsection{Anytime-valid monitoring and many alarms}
Closed-view detectors can be anytime-valid (confidence sequences, e-processes)~\cite{howard2021confseq,ramdas2024evalues}.
The default Ville/Markov rejection threshold $1/\alpha$ for e-processes is often conservative; recent work gives sharper thresholds under mild, empirically checkable shape assumptions on the e-value distribution~\cite{blierwong2026thresholds}.
But if many detectors exist, naive ``raise any alarm'' logic incurs multiple-testing inflation.
E-values provide dependence-robust FDR control (e-BH)~\cite{wang2022ebh}.
When alarms are raised at global stopping times, one should use stopped e-BH variants~\cite{wang2025sebH}.

\begin{figure}[t]
  \centering
  \fbox{\parbox{0.92\linewidth}{\centering\textbf{Figure omitted in source-only archive.}\\(Plot/diagram and scripts available on request.)}}
  \caption{Example of an anytime-valid monitoring boundary for adaptive/streaming settings.}
\end{figure}


\section{CPPC manifests}
\label{sec:cppc}

We propose \textbf{Control-Plane Privacy Cards (CPPCs)}: short, machine-checkable manifests describing control-plane witness channels and their budgets.
A CPPC is not a proof; it is a contract that enables third-party auditing and comparison.

\paragraph{What a CPPC declares.}
At minimum:
(i) which control-plane signals are released (weights, beacons, alerts),
(ii) cadence and granularity,
(iii) accounting model (e.g., max-divergence DP, RDP, or tradeoff-curve notions such as $f$-DP/GDP~\cite{dong2022gdp}) and composition regime,
(iv) refresh/expiration policies,
(v) how monitoring/alarms are handled (closed view vs public),
(vi) the \emph{unit of privacy} / adjacency notion and any interpretability guidance for the chosen budget representation (e.g., as emphasized by NIST's guidance for evaluating DP guarantees~\cite{nist_sp800_226}).

\begin{figure}[t]
  \centering
  \fbox{\parbox{0.92\linewidth}{\centering\textbf{Figure omitted in source-only archive.}\\(Plot/diagram and scripts available on request.)}}
  \caption{A ``compiler'' view of CPPCs: selection + releases + monitoring are compiled into a single budgeted interface.}
\end{figure}

\begin{figure}[t]
  \centering
  \fbox{\parbox{0.92\linewidth}{\centering\textbf{Figure omitted in source-only archive.}\\(Plot/diagram and scripts available on request.)}}
  \caption{Example CPPC cards.}
\end{figure}

\paragraph{Publication and integrity.}
In practice, CPPCs should be signed and published via an append-only transparency mechanism (Merkle-tree log),
analogous to Certificate Transparency~\cite{rfc9162_ct}.
This is also in the spirit of ``documentation artifacts'' such as model cards and datasheets~\cite{mitchell2019modelcards,gebru2021datasheets},
and can reuse supply-chain primitives such as SBOM formats (e.g., SPDX or CycloneDX)~\cite{spdx_overview,cyclonedx_overview}.

\paragraph{Verification sketch.}
A minimal deployment can be lightweight:
(i) the operator signs a CPPC (and version bumps) and submits it to a public log;
(ii) each public release channel (weights, beacons, alerts) includes a short log pointer (e.g., inclusion proof or log checkpoint);
(iii) auditors verify the signature and log inclusion, then check \emph{consistency} between CPPC declarations and observable releases:
cadence/granularity match; budget beacons are monotone and compose according to the declared accountant; and any public alerting behavior
matches the declared ``closed view vs public'' handling.
This does not reveal raw telemetry, but it makes covert, out-of-contract control-plane channels harder to hide.

\begin{table}[t]
\centering
\begin{tabular}{@{}p{0.45\linewidth}p{0.45\linewidth}@{}}
\toprule
\textbf{Auditor can verify from public artifacts} & \textbf{Typically not verifiable without trust/attestation} \\
\midrule
CPPC signature validity and log inclusion (inclusion/consistency proofs)~\cite{rfc9162_ct,tomescu2019transparencylogs,newman2022sigstore} & Correctness of internal telemetry, sampling, and detector implementations \\
Release cadence and granularity match CPPC (timestamps, bucket sizes) & Whether additional closed-view channels exist beyond those declared \\
Budget beacon monotonicity and declared composition schedule (odometer/filter) & The operator's internal ``stop time'' logic (e.g., whether early alarms were suppressed) \\
Public alert policy matches CPPC (fixed schedule vs data-dependent, use of randomized response) & Confidential incident details and on-call runbooks \\
Artifacts referenced by the CPPC (SBOM-like metadata, versioning) are present and consistent~\cite{spdx_spec301,slsa_spec_v10,in_toto2019} & Whether budgeted releases were computed from the claimed inputs (needs reproducibility or trusted execution) \\
\bottomrule
\end{tabular}
\caption{A minimal checklist for ``budget-consistency'' auditing. The goal is not to prove security from public data, but to make it difficult to silently change witness channels, cadences, or budgets without leaving detectable inconsistencies.}
\label{tab:auditor-checklist}
\end{table}


\section{Related work}
\paragraph{Accounting.}
Differential privacy provides the canonical stability language for releasing summaries~\cite{dworkroth2014foundations}.
Odometers and filters address adaptive-parameter composition~\cite{rogers2016odometers}, with practical RDP-oriented variants~\cite{lecuuyer2021practical,mironov2017rdp}.
Interactive deployments motivate concurrent composition results~\cite{vadhanwang2021concurrentdp}.
Privacy amplification by subsampling supports the ``closed-view'' intuition that partial/public views can be strictly safer~\cite{balle2018subsampling}.
Tradeoff-curve notions such as $f$-DP (and its Gaussian specialization, GDP) preserve the testing interpretation and enable sharp composition reasoning~\cite{dong2022gdp}.
Recent work studies when fully-adaptive $f$-DP filters are valid and gives approximate GDP filters for practical regimes~\cite{tran2026fdpfilters}.
Natural (profile-level) privacy filters promise sharper accounting but are \emph{not} always ``free'' under adaptive composition, even for common families of mechanisms~\cite{regehr2026naturalfilters}.
``Accuracy-first'' and ex-post privacy perspectives offer another interface for adaptively choosing privacy parameters while preserving post-processing immunity in the accounting model~\cite{raisa2025accuracyfirstrdp}.


\paragraph{Always-valid monitoring and multiple testing.}
Confidence sequences and e-values support anytime-valid inference in streaming systems~\cite{howard2021confseq,ramdas2024evalues}.
E-values yield dependence-robust FDR control via e-BH~\cite{wang2022ebh}; stopped e-BH variants address global stopping-time alerts~\cite{wang2025sebH}.
Recent work applies these tools to \emph{audit} privacy guarantees from sampled outputs, including sequential auditors for $(\varepsilon,\delta)$-DP and for $f$-DP~\cite{gonzalez2025sequentiallyauditdp,kutta2026sequentialauditfdp,mahloujifar2025auditfdpone}.


\paragraph{Private telemetry systems.}
Prio and RAPPOR demonstrate scalable private aggregation and telemetry in production-like settings~\cite{corrigangibbs2017prio,erlingsson2014rappor}.

\paragraph{Transparency and ``cards.''}
Model cards and datasheets helped normalize lightweight, comparable documentation for models/datasets~\cite{mitchell2019modelcards,gebru2021datasheets}.
Supply-chain work provides analogous machine-checkable manifests (SBOM) and provenance: SPDX and CycloneDX, SLSA (v1.0 is marked as ``retired as of 2026), and in-toto~\cite{spdx_spec301,slsa_spec_v10,in_toto2019}.
Append-only transparency logs (Certificate Transparency) provide a simple publication primitive~\cite{rfc9162_ct}, while Sigstore aims to make signing/provenance usable at scale~\cite{newman2022sigstore}.
NIST's Secure Software Development Framework (SSDF) is a recent baseline for secure development practice~\cite{nist_ssdf_800_218}.
CPPCs adapt these documentation/provenance ideas specifically to anonymity control planes.




\section{Takeaways}
\begin{itemize}[leftmargin=*]
\item \textbf{Monitoring is a witness stream:} alarms, trigger times, and public health beacons can act as sequential identification oracles and must be priced.
\item \textbf{Prefer closed view:} keep rich telemetry and detectors inside a measurement firewall; release only coarse, precommitted summaries governed by an accountant.
\item \textbf{Account adaptivity explicitly:} use odometers/filters (or $f$-DP/GDP tradeoff accounting) when thresholds, cadences, or detector libraries change over time.
\item \textbf{Make the contract machine-checkable:} CPPCs and transparency logging make it harder to silently change witness channels, cadences, or budgets without leaving a detectable inconsistency.
\end{itemize}

\begin{thebibliography}{99}\small

\bibitem{dworkroth2014foundations}
Cynthia Dwork and Aaron Roth.
\newblock The Algorithmic Foundations of Differential Privacy.
\newblock \emph{Foundations and Trends in Theoretical Computer Science}, 9(3--4):211--407, 2014.

\bibitem{mironov2017rdp}
Ilya Mironov.
\newblock Rényi Differential Privacy.
\newblock In \emph{CSF}, 2017.\; \texttt{arXiv:1702.07476}.

\bibitem{balle2020rdp_testing}
Borja Balle.
\newblock Hypothesis Testing Interpretations and Rényi Differential Privacy.
\newblock In \emph{ICML (PMLR 108)}, 2020.

\bibitem{dwork2010continual}
Cynthia Dwork, Moni Naor, Toniann Pitassi, and Guy Rothblum.
\newblock Differential Privacy Under Continual Observation.
\newblock In \emph{STOC}, 2010.

\bibitem{chan2010continual}
T.-H. Hubert Chan, Elaine Shi, and Dawn Song.
\newblock Private and Continual Release of Statistics.
\newblock \texttt{ePrint:2010/076}. (See also TCC/ICALP-era versions.)

\bibitem{jain2023priceContinual}
Palak Jain, Sofya Raskhodnikova, Satchit Sivakumar, and Adam Smith.
\newblock The Price of Differential Privacy under Continual Observation.
\newblock In \emph{ICML (PMLR 202)}, 2023.

\bibitem{cohen2024continual_lowerbounds}
Edith Cohen, Xin Lyu, Jelani Nelson, Tamás Sarlós, and Uri Stemmer.
\newblock Lower Bounds for Differential Privacy Under Continual Observation and Online Threshold Queries.
\newblock In \emph{ICML (PMLR 247)}, 2024.

\bibitem{rogers2016odometers}
Ryan Rogers, Aaron Roth, Jonathan Ullman, and Salil Vadhan.
\newblock Privacy Odometers and Filters: Pay-as-you-Go Composition.
\newblock In \emph{NeurIPS}, 2016.\; \texttt{arXiv:1605.08294}.

\bibitem{vadhanwang2021concurrentdp}
Salil Vadhan and Tianhao Wang.
\newblock Concurrent Composition of Differential Privacy.
\newblock In \emph{TCC}, 2021.

\bibitem{lecuuyer2021practical}
Mathias Lécuyer et al.
\newblock Practical Privacy Filters and Odometers with Rényi Differential Privacy (and Applications to DP Deep Learning).
\newblock \texttt{arXiv:2103.01379}, 2021.

\bibitem{erlingsson2014rappor}
Úlfar Erlingsson, Vasyl Pihur, and Aleksandra Korolova.
\newblock RAPPOR: Randomized Aggregatable Privacy-Preserving Ordinal Response.
\newblock In \emph{CCS}, 2014.\; \texttt{arXiv:1407.6981}.

\bibitem{corrigangibbs2017prio}
Henry Corrigan-Gibbs and Dan Boneh.
\newblock Prio: Private, Robust, and Scalable Computation of Aggregate Statistics.
\newblock In \emph{NSDI}, 2017.\; \texttt{arXiv:1703.06255}.

\bibitem{howard2021confseq}
Steven R. Howard et al.
\newblock Time-uniform, nonparametric, nonasymptotic confidence sequences.
\newblock \emph{Annals of Statistics}, 49(2), 2021.\; \texttt{arXiv:1810.08240}.

\bibitem{ramdas2024evalues}
Aaditya Ramdas.
\newblock Hypothesis Testing with E-values.
\newblock \texttt{arXiv:2410.23614}, 2024.

\bibitem{wang2022ebh}
Ruodu Wang.
\newblock False discovery rate control with e-values.
\newblock \emph{JRSS~B}, 84(3):822--852, 2022.

\bibitem{wang2025sebH}
Hongjian Wang, Sanjit Dandapanthula, and Aaditya Ramdas.
\newblock Anytime-valid FDR control with the stopped e-BH procedure.
\newblock \texttt{arXiv:2502.08539}, 2025.

\bibitem{blierwong2026thresholds}
Christopher Blier-Wong and Ruodu Wang.
\newblock Improved thresholds for e-values.
\newblock Preprint / early 2026 versions available online.

\bibitem{gebru2021datasheets}
Timnit Gebru et al.
\newblock Datasheets for Datasets.
\newblock \emph{Communications of the ACM}, 64(12), 2021.\; (also \texttt{arXiv:1803.09010}).

\bibitem{mitchell2019modelcards}
Margaret Mitchell et al.
\newblock Model Cards for Model Reporting.
\newblock In \emph{FAT*}, 2019.\; \texttt{arXiv:1810.03993}.

\bibitem{newman2022sigstore}
Zachary Newman et al.
\newblock Sigstore: Software Signing for Everybody.
\newblock In \emph{CCS}, 2022.

\bibitem{in_toto2019}
Santiago Torres-Arias et al.
\newblock in-toto: Providing farm-to-table guarantees for bits and bytes.
\newblock In \emph{USENIX Security}, 2019.

\bibitem{slsa_spec_v10}
SLSA.
\newblock SLSA specification v1.0 (Status: Retired).
\newblock \url{https://slsa.dev/spec/v1.0/}.

\bibitem{cyclonedx_overview}
OWASP CycloneDX.
\newblock Specification overview.
\newblock \url{https://cyclonedx.org/specification/overview/}.

\bibitem{spdx_overview}
SPDX.
\newblock SPDX specification overview.
\newblock \url{https://spdx.github.io/spdx-spec/v3.0.1/}.

\bibitem{spdx_spec301}
SPDX.
\newblock SPDX Specification v3.0.1.
\newblock \url{https://spdx.dev/wp-content/uploads/sites/31/2024/12/SPDX-3.0.1-1.pdf}.


\bibitem{series:schedcompiler}
Anonymous.
\newblock \emph{Scheduling and Compiler Techniques for Metadata-Hiding Overlay Lookups}.
\newblock Series paper (published), 2026.
\newblock Wiki: \texttt{[[2026.01.22 - Mathematics: Scheduling and Compiler Techniques for Metadata-Hiding Overlay Lookups]]}.

\bibitem{series:stoptimeaddendum}
Anonymous.
\newblock \emph{Stop-Time Padding Addendum (unique-only): Max-Mixture Separations and Tail-Sign Deadline Normal Forms}.
\newblock Series note (published), 2026.
\newblock Wiki: \texttt{[[2026.01.26 - Mathematics: Stop-Time Padding Addendum, Max-Mixture Separations and Tail-Sign Deadline Normal Forms]]}.

\bibitem{series:state1}
Anonymous.
\newblock Stateful Anonymity and Control-Plane Leakage: Selection Privacy, Amplification, and Accounting.
\newblock Companion preprint in this archive (AnonDHT~State~1).

\bibitem{balle2018subsampling}
Borja Balle and Gilles Barthe and Marco Gaboardi.
\newblock Privacy Amplification by Subsampling: Tight Analyses via Couplings and Divergences.
\newblock Advances in Neural Information Processing Systems (NeurIPS), 2018.
\newblock NeurIPS 2018; also available as arXiv:1807.01647

\bibitem{nist_ssdf_800_218}
{National Institute of Standards and Technology}.
\newblock Secure Software Development Framework (SSDF) Version 1.1: Recommendations for Mitigating the Risk of Software Vulnerabilities.
\newblock NIST Special Publication 800-218, 2022.
\newblock \url{https://csrc.nist.gov/pubs/sp/800/218/final}

\bibitem{nist_sp800_226}
Joseph P. Near and David Darais and Naomi Lefkovitz and Gary S. Howarth.
\newblock Guidelines for Evaluating Differential Privacy Guarantees.
\newblock NIST Special Publication 800-226, 2025.
\newblock DOI: 10.6028/NIST.SP.800-226.
\newblock Final publication (March 2025) with PDF available from the NIST landing page.
\newblock \url{https://csrc.nist.gov/pubs/sp/800/226/final}

\bibitem{xie2025dpcusum}
Liyan Xie and Ruizhi Zhang.
\newblock Sequential Change Detection with Differential Privacy.
\newblock 2025.
\newblock DOI: 10.48550/arXiv.2509.02768.
\newblock arXiv:2509.02768
\newblock \url{https://arxiv.org/abs/2509.02768}

\bibitem{gonzalez2025sequentiallyauditdp}
Tom{\'a}s Gonz{\'a}lez and Mateo Dulce-Rubio and Aaditya Ramdas and M{\'o}nica Ribero.
\newblock Sequentially Auditing Differential Privacy.
\newblock 2025.
\newblock arXiv:2509.07055; also presented at NeurIPS 2025.
\newblock \url{https://arxiv.org/abs/2509.07055}

\bibitem{kutta2026sequentialauditfdp}
Tim Kutta and Martin Dunsche and Yu Wei and Vassilis Zikas.
\newblock Sequential Auditing for $f$-Differential Privacy.
\newblock 2026.
\newblock arXiv:2602.06518

\bibitem{mahloujifar2025auditfdpone}
Saeed Mahloujifar and Luca Melis and Kamalika Chaudhuri.
\newblock Auditing $f$-Differential Privacy in One Run.
\newblock Proceedings of the 42nd International Conference on Machine Learning (ICML), 267, 42615--42641, 2025.
\newblock Also available as arXiv:2410.22235 and via OpenReview.
\newblock \url{https://proceedings.mlr.press/v267/mahloujifar25a.html}

\bibitem{dong2022gdp}
Jinshuo Dong and Aaron Roth and Weijie J. Su.
\newblock Gaussian Differential Privacy.
\newblock Journal of the Royal Statistical Society: Series B (Statistical Methodology), 84(1), 3--37, 2022.
\newblock DOI: 10.1111/rssb.12455.

\bibitem{tran2026fdpfilters}
Long Tran and Antti Koskela and Ossi R{"a}is{"a} and Antti Honkela.
\newblock $f$-Differential Privacy Filters: Validity and Approximate Solutions.
\newblock 2026.
\newblock DOI: 10.48550/arXiv.2602.06756.
\newblock arXiv:2602.06756
\newblock \url{https://arxiv.org/abs/2602.06756}

\bibitem{regehr2026naturalfilters}
Matthew Regehr and Bingshan Hu and Ethan Leeman and Pasin Manurangsi and Pierre Tholoniat and Mathias L{\'e}cuyer.
\newblock Natural Privacy Filters Are Not Always Free: A Characterization of Free Natural Filters.
\newblock 2026.
\newblock DOI: 10.48550/arXiv.2602.15815.
\newblock arXiv:2602.15815
\newblock \url{https://arxiv.org/abs/2602.15815}

\bibitem{raisa2025accuracyfirstrdp}
Ossi R{\"a}is{\"a} and Antti Koskela and Antti Honkela.
\newblock Accuracy-First R{\'e}nyi Differential Privacy and Post-Processing Immunity.
\newblock 2025.
\newblock DOI: 10.48550/arXiv.2509.22213.
\newblock arXiv:2509.22213
\newblock \url{https://arxiv.org/abs/2509.22213}

\bibitem{rfc9162_ct}
Ben Laurie and Adam Langley and Emilia Kasper and Eran Messeri.
\newblock Certificate Transparency Version 2.0.
\newblock 2021.
\newblock RFC 9162.
\newblock \url{https://www.rfc-editor.org/rfc/rfc9162.html}

\bibitem{tomescu2019transparencylogs}
Alin Tomescu and Vivek Bhupatiraju and Dimitrios Papadopoulos and Charalampos Papamanthou and Nikos Triandopoulos and Srinivas Devadas.
\newblock Transparency Logs via Append-Only Authenticated Dictionaries.
\newblock Proceedings of the 2019 {ACM} SIGSAC Conference on Computer and Communications Security (CCS '19), 1299--1316, 2019.
\newblock DOI: 10.1145/3319535.3345652.
\newblock General-purpose transparency log construction and security formalization.

\bibitem{warner1965randomizedresponse}
Stanley L. Warner.
\newblock Randomized Response: A Survey Technique for Eliminating Evasive Answer Bias.
\newblock Journal of the American Statistical Association, 60(309), 63--69, 1965.
\newblock DOI: 10.1080/01621459.1965.10480775.

\bibitem{pedersen1991distributedprovers}
Torben P. Pedersen.
\newblock Distributed Provers with Applications to Undeniable Signatures.
\newblock Advances in Cryptology --- {EUROCRYPT} '91, 547, 221--242, 1991.
\newblock DOI: 10.1007/3-540-46416-6\_20.

\bibitem{series:artifactpolicy}
Anonymous.
\newblock Artifact policy and supporting materials for the anonymity/AnonDHT series.
\newblock 2026.
\newblock See synthesis paper: Artifact policy (this archive). Artifacts, tooling, certificates, and scripts are available on request as a single archive.
\bibitem{series:contractobjects}
Anonymous.
\newblock Contract Objects for Anonymous-DHT Privacy Budgets and Claims.
\newblock Series synthesis note (Synthesis~0), 2026. (This archive.)

\bibitem{series:abomoinl}
Anonymous.
\newblock ABOM/OINL: Audit BOM and Outcome Identification Noise Label for Anonymous DHT Deployments.
\newblock Anonymity series Paper~C, 2026. (This archive.)

\bibitem{series:state3addendum}
Anonymous.
\newblock Closed-View Auditing and CPPCs (Addendum): Multi-Stream Monitoring Notes and Schema Excerpts.
\newblock AnonDHT State series Paper~3A, 2026. (This archive.)

\bibitem{series:traceifaces}
Anonymous.
\newblock Trace Interfaces for Anonymous {DHT}s: Observation Tiers, Committee-Contact Traces, and Exposure Normal Forms.
\newblock Series synthesis note (Synthesis~3), 2026. (This archive.)

\bibitem{series:threatwindows}
Anonymous.
\newblock Threat Windows for Anonymous {DHT}s: What a Budget Covers and For How Long.
\newblock Series synthesis note (Synthesis~4), 2026. (This archive.)


\bibitem{series:notation}
Anonymous.
\newblock Notation and Minimal Model for the One-True-Archive.
\newblock Series synthesis note (Synthesis~8), 2026. (This archive.)

\end{thebibliography}

\end{document}